% ============================================================
% Capítulo: Resultados e Discussão
% ============================================================

\chapter{Resultados e Discussão}
\label{cap:resultados}

Este capítulo apresenta e discute os resultados obtidos pelos modelos de \textit{parsing} de
dependências desenvolvidos neste trabalho, avaliados no conjunto de teste do corpus
Porttinari. A análise está organizada em cinco eixos principais: desempenho
geral comparado ao PortParser; impacto da arquitetura do cabeçote de predição; efeito do MTL;
análise de desempenho por rótulo de dependência; e dinâmica de instabilidade durante o
treinamento. As definições completas dos rótulos de dependência citados ao longo do capítulo
encontram-se no Apêndice~\ref{cap:classe_deprel}. Os resultados detalhados da busca de
hiperparâmetros e as curvas de treinamento por época de cada modelo estão disponíveis,
respectivamente, nos Apêndices~\ref{cap:resultados_hiperparametros}
e~\ref{cap:modelos_avaliados}.

% ─────────────────────────────────────────────────────────────────────────────
\section{Desempenho Geral}
\label{sec:desempenho_geral}
% ─────────────────────────────────────────────────────────────────────────────

A Tabela~\ref{tab:comparation} e a Figura~\ref{fig:resultados_geral} apresentam os resultados
de todos os modelos avaliados em termos de \textit{Unlabeled Attachment Score} (UAS),
\textit{Labeled Attachment Score} (LAS) e acurácia de UPOS. O melhor resultado obtido neste
trabalho foi alcançado pelo modelo \textbf{BERTimbau Large com cabeçote linear e aprendizado
multi-tarefa} (\textit{BtL-Lin-MTL}), com UAS de 94,99\% e LAS de 93,84\%.

\begin{figure}[ht]
  \centering
  \includegraphics[width=\textwidth]{Dissertação/Imagens/resultados_geral.pdf}
  \caption{Comparação de desempenho (UAS e LAS) de todos os modelos avaliados, ordenados por UAS
           decrescente. As linhas tracejadas indicam o desempenho do PortParser como referência
           (\textbf{UAS\,=\,96,08\%}, \textbf{LAS\,=\,94,61\%}).
           cU\,=\,com UPOS (MTL); sU\,=\,sem UPOS (ablação).}
  \label{fig:resultados_geral}
\end{figure}

Em relação ao \textbf{PortParser}~\citep{lopes2024towards}, modelo de referência para o português
brasileiro, o melhor resultado deste trabalho apresenta uma diferença de \textbf{1,09 pontos
percentuais} em UAS e \textbf{0,77 pontos percentuais} em LAS. A principal diferença arquitetural
que explica essa lacuna é a \textbf{ausência de uma camada BiLSTM bidirecional} no ParseH2IA:
enquanto o PortParser insere uma camada BiLSTM entre o encoder e os cabeçotes de predição para
contextualização adicional e realiza ajuste fino completo, o ParseH2IA conecta o encoder
diretamente aos cabeçotes com ajuste fino completo de todos os parâmetros, sem camada
intermediária. Essa arquitetura mais simples permite uma análise controlada da contribuição
do cabeçote de predição isolado da influência do BiLSTM.

Nota-se que a lacuna de desempenho é consideravelmente pequena dada a supressão do BiLSTM:
o BtL-Lin-MTL aproxima-se do PortParser a menos de 1,1 pontos percentuais em UAS, sugerindo
que o BERTimbau Large com ajuste fino completo produz representações suficientemente ricas para
suportar o \textit{parsing} diretamente, sem a camada de contextualização adicional.

Quanto à \textbf{acurácia de UPOS}, todos os modelos MTL superam o PortParser
(99,09\%), com valores entre 98,81\% (JabuticaBERT Linear) e 99,34\% (BERTimbau Large
Biaffine). Esse resultado indica que a tarefa de classificação morfológica está próxima do
limite superior de desempenho para o Porttinari, independentemente da arquitetura do cabeçote.

% ─────────────────────────────────────────────────────────────────────────────
\section{Comparação entre Encoders}
\label{sec:encoders}
% ─────────────────────────────────────────────────────────────────────────────

A Figura~\ref{fig:curvas_enc} mostra as curvas de aprendizado (UAS no conjunto de teste por
época) para os modelos com cabeçote linear e MTL de cada encoder. A hierarquia de desempenho
final é consistente ao longo de praticamente todo o treinamento:

\[
  \text{BERTimbau Large} > \text{BERTimbau Base} > \text{mBERT} > \text{JabuticaBERT}
\]

\begin{figure}[ht]
  \centering
  \includegraphics[width=0.92\textwidth]{Dissertação/Imagens/curvas_treinamento_enc.pdf}
  \caption{Curvas de UAS por época para os modelos Linear MTL de cada encoder. O eixo vertical
           indica o UAS (\%) no conjunto de teste ao final de cada época.}
  \label{fig:curvas_enc}
\end{figure}

O \textbf{BERTimbau Large} lidera em todas as configurações, confirmando que modelos pré-treinados
em português com maior capacidade oferecem representações mais adequadas para a tarefa. O
\textbf{BERTimbau Base}, treinado no mesmo corpus mas com menor capacidade, apresenta desempenho
próximo ao Large, com diferença de apenas 1,25 pontos percentuais em UAS no melhor caso
(94,99\% vs.\ 93,74\%).

O \textbf{mBERT}, apesar de ser multilíngue, supera o JabuticaBERT em todas as métricas, exceto
na variante Biaffine sem UPOS. Isso é notável, pois o JabuticaBERT é um modelo monolíngue para o
português. A explicação mais provável está no tamanho e arquitetura do JabuticaBERT: trata-se de
um modelo \textit{Base} treinado com técnicas de aprendizado modernas (\textit{continual
pretraining}, vocabulário expandido), o que pode gerar uma distribuição de representações
subótima para a interface direta com os cabeçotes de predição. Além
disso, o JabuticaBERT utiliza um vocabulário maior e tokenização diferente do BERTimbau,
gerando potencialmente mais subpalavras por token e introduzindo ruído na etapa de extração
do primeiro subtoken por palavra.

As curvas de aprendizado revelam ainda um padrão de \textbf{convergência em duas fases} para os
modelos com mais épocas: um crescimento rápido nas primeiras 20 épocas (fase de alta taxa de
aprendizado, $\eta = 10^{-3}$), seguido de uma fase de refinamento mais lenta. As curvas
completas por época, incluindo \textit{training loss} e \textit{validation loss}, estão
disponíveis no Apêndice~\ref{cap:modelos_avaliados} e corroboram que todos os modelos
convergiram sem sinais de sobreajuste grave na fase inicial, com a \textit{validation loss}
estabilizando antes do fim do treinamento. Esse comportamento é coerente com o esquema de
taxa de aprendizado em dois estágios adotado durante o treinamento, cujos hiperparâmetros
ótimos por encoder foram determinados pela busca apresentada no
Apêndice~\ref{cap:resultados_hiperparametros}.

% ─────────────────────────────────────────────────────────────────────────────
\section{Impacto da Arquitetura do Cabeçote}
\label{sec:arquitetura}
% ─────────────────────────────────────────────────────────────────────────────

\subsection{Linear versus Biaffine}

O resultado mais surpreendente deste trabalho é a \textbf{superioridade consistente dos cabeçotes
lineares sobre os biaffines} em UAS e LAS (Figura~\ref{fig:linear_biaffine}). Em três dos quatro
encoders (BtL, BtB e mBERT), o cabeçote linear supera o biaffine em UAS por uma margem
substancial: 3,47 pontos para BERTimbau Large (94,99\% vs.\ 91,52\%), 3,90 pontos para
BERTimbau Base (93,74\% vs.\ 89,84\%) e 5,78 pontos para mBERT (92,70\% vs.\ 86,92\%). A
única exceção é o JabuticaBERT, onde o Biaffine supera levemente o Linear (88,94\% vs.\ 87,78\%
em UAS para a variante MTL, e 90,04\% vs.\ 87,00\% sem UPOS).

Vale observar que essa diferença não é artefato da escolha de hiperparâmetros: como mostra o
Apêndice~\ref{cap:resultados_hiperparametros}, a busca bayesiana de hiperparâmetros foi
conduzida \textbf{independentemente} para cada combinação de encoder e arquitetura, e os
hiperparâmetros ótimos do biaffine foram aplicados no treinamento final. Para o BERTimbau
Large, por exemplo, os melhores hiperparâmetros do biaffine (Trial~0: LR\,=\,$3{,}0\times
10^{-5}$, \textit{weight decay}\,=\,0,116, \textit{warmup ratio}\,=\,0,417) produziram
LAS médio de 0,887 na validação cruzada — substancialmente inferior ao LAS médio do linear
(0,926, Trial~2). A diferença persiste no conjunto de teste final, confirmando que a superioridade
do linear é sistemática e não decorre de sub-otimização do biaffine.

\begin{figure}[ht]
  \centering
  \includegraphics[width=0.92\textwidth]{Dissertação/Imagens/linear_vs_biaffine.pdf}
  \caption{UAS (\%) por variante de arquitetura (Lin/Biaff × cU/sU) e encoder. A superioridade
           do cabeçote linear sobre o biaffine é consistente em três dos quatro encoders.}
  \label{fig:linear_biaffine}
\end{figure}

Na literatura, o cabeçote biaffine é amplamente reconhecido como superior para o
\textit{parsing} de dependências~\citep{dozat2016deep}, o que torna este resultado inesperado.
Quatro hipóteses, complementares entre si, explicam esse comportamento no contexto específico
deste trabalho:

\begin{enumerate}
  \item \textbf{Codificação sintática implícita nas representações BERT}: Estudos de
        \textit{probing}~\citep{tenney2019bert,hewitt2019structural,jawahar2019does} demonstraram
        que modelos de linguagem pré-treinados como o BERT codificam, de forma implícita,
        informação sintática substancial em seus vetores de saída — incluindo relações de
        dependência, distâncias na árvore sintática e papéis gramaticais. Mesmo após o ajuste
        fino completo, essa organização sintática do espaço de representações tende a ser
        preservada~\citep{hewitt2019structural}. Nesse cenário, um cabeçote linear atua como
        uma \textit{sonda linear} (\textit{linear probe}) direta sobre a estrutura já organizada
        pelo encoder, aprendendo de forma eficiente e estável. O cabeçote biaffine, ao introduzir
        transformações MLP e bilineares adicionais, adiciona graus de liberdade que, em corpus
        pequeno, podem não convergir para uma solução melhor do que a leitura direta pelo linear.
        Essa hipótese é corroborada pelo fato de que a vantagem do linear é mais pronunciada
        nos encoders de maior capacidade (BERTimbau Large e Base), que tendem a codificar mais
        ricamente a estrutura sintática em suas representações.

  \item \textbf{Corpus limitado e sobreajuste}: O cabeçote biaffine introduz substancialmente
        mais parâmetros treináveis: quatro MLPs (Linear\,+\,ELU\,+\,LayerNorm\,+\,Dropout)
        com $d_\text{arc}=500$ e $d_\text{rel}=100$, mais duas matrizes bilineares. Com apenas
        5.893 sentenças de treinamento no Porttinari, essa capacidade adicional aumenta o risco
        de sobreajuste às particularidades do corpus, enquanto o cabeçote linear — com suas
        três projeções independentes — atua como regularização implícita.

  \item \textbf{Instabilidade de treinamento}: Como será discutido na Seção~\ref{sec:instabilidade},
        o cabeçote biaffine apresenta maior número de eventos de \textit{forgetting} durante o
        treinamento. As curvas detalhadas no Apêndice~\ref{cap:modelos_avaliados} confirmam que
        os modelos biaffines exibem oscilações mais acentuadas entre \textit{validation loss} e
        UAS ao longo das épocas, sugerindo que sua otimização é mais instável em corpus pequeno,
        o que prejudica o desempenho final.
\end{enumerate}

A hipótese da codificação sintática implícita é particularmente relevante porque inverte a
lógica da escolha arquitetural usual: ao invés de projetar um cabeçote mais expressivo para
capturar relações complexas, o encoder com ajuste fino já capturou essas relações — e o papel
do cabeçote é simplesmente lê-las, tarefa para a qual a simplicidade do linear é uma vantagem,
não uma limitação, especialmente com corpus de treinamento restrito.

\subsection{Impacto do Multi-Task Learning}

A Figura~\ref{fig:impacto_mtl} mostra a diferença de UAS e LAS entre os modelos com e sem a
tarefa auxiliar de UPOS (MTL versus ablação). O impacto geral é \textbf{pequeno e assimétrico}:
o MTL beneficia consistentemente os modelos lineares (ganho médio de +0,33 pp em UAS para o
BERTimbau), mas prejudica os modelos biaffines (perda média de $-$1,10 pp em UAS), com a exceção
do JabuticaBERT Biaffine.

\begin{figure}[ht]
  \centering
  \includegraphics[width=0.92\textwidth]{Dissertação/Imagens/impacto_mtl.pdf}
  \caption{Diferença de UAS e LAS (em pontos percentuais) entre modelos MTL (com UPOS) e de
           ablação (sem UPOS), por variante e encoder. Valores positivos indicam que o MTL
           beneficia o \textit{parsing}; negativos indicam prejuízo.}
  \label{fig:impacto_mtl}
\end{figure}

Esse resultado sugere que a tarefa de UPOS — já próxima do teto de desempenho ($>$99\%) — adiciona
pouca informação de gradiente útil ao ajuste fino do encoder sobre o corpus Porttinari. O ganho
marginal nos modelos lineares pode decorrer de uma leve regularização implícita introduzida pela
tarefa auxiliar, enquanto o prejuízo nos modelos biaffines pode resultar de um conflito de
gradientes: as tarefas de HEAD e UPOS competem pelos parâmetros das MLPs com sinais de grandezas
muito diferentes, prejudicando a convergência do mecanismo biaffine.

% ─────────────────────────────────────────────────────────────────────────────
\section{Análise por Rótulo de Dependência}
\label{sec:por_rotulo}
% ─────────────────────────────────────────────────────────────────────────────

A análise por rótulo DEPREL (Tabelas~\ref{tab:acc-mbert}--\ref{tab:acc-jab} no
Apêndice~\ref{ap:tabelas}) revela uma variação expressiva de desempenho entre rótulos. As
definições de cada rótulo citado nesta seção encontram-se no Apêndice~\ref{cap:classe_deprel}.

\subsection{Rótulos de Alta Acurácia}

Os rótulos \texttt{punct} (0,9982--0,9991), \texttt{case} (0,9956--0,9982) e \texttt{cop}
(0,9849--0,9983) apresentam alta acurácia consistente em todos os modelos. Esses rótulos
correspondem a palavras com forte restrição gramatical: pontuação, preposições/posposições e
verbos de ligação têm pouca ambiguidade de dependência e são frequentes no Porttinari
(para suas definições, ver Apêndice~\ref{cap:classe_deprel}). A alta frequência de
ocorrência fornece mais exemplos de treinamento, enquanto a baixa ambiguidade resulta em
padrões altamente previsíveis. Esses são também os rótulos mais facilmente extraíveis por um
cabeçote linear, pois a informação sintática correspondente tende a estar fortemente codificada
nas representações do BERT — reforçando a hipótese de leitura direta discutida anteriormente.

\subsection{Rótulos Difíceis}

A Figura~\ref{fig:labels_dificeis} apresenta os rótulos com menor acurácia nos quatro melhores
modelos. Os rótulos \texttt{orphan}, \texttt{dislocated}, \texttt{vocative} e \texttt{parataxis}
são consistentemente os mais desafiadores (ver Apêndice~\ref{cap:classe_deprel} para
definições detalhadas).

\begin{figure}[ht]
  \centering
  \includegraphics[width=0.92\textwidth]{Dissertação/Imagens/labels_dificeis.pdf}
  \caption{Acurácia (\%) nos rótulos DEPREL mais difíceis para os quatro melhores modelos.
           \texttt{orphan} e \texttt{dislocated} apresentam desempenho inferior a 50\% na
           maioria dos modelos.}
  \label{fig:labels_dificeis}
\end{figure}

\paragraph{orphan e dislocated.} Com acurácia abaixo de 50\% na maioria dos modelos, esses
rótulos são raros no Porttinari (14 e 17 ocorrências no conjunto de teste, respectivamente) e
semanticamente complexos: \texttt{orphan} marca dependentes de um núcleo elidido em estruturas
de elipse, enquanto \texttt{dislocated} é usado em deslocamentos sintáticos
(ver Apêndice~\ref{cap:classe_deprel}). A raridade extrema impede o aprendizado de padrões
robustos, e a natureza não-local dessas construções dificulta sua codificação implícita pelo
BERT, que opera com viés de proximidade na atenção.

\paragraph{parataxis e discourse.} Esses rótulos codificam relações semanticamente difusas,
com acurácia variando de 62,75\% a 78,12\% para \texttt{parataxis} e de 56,86\% a 74,51\% para
\texttt{discourse}. Ambos são frequentemente confundidos entre si e com outros rótulos de alta
abstração (e.g., \texttt{advcl}, \texttt{ccomp}). Trata-se de relações que dependem de
conhecimento discursivo e pragmático que o BERT capta apenas parcialmente via \textit{masked
language modeling}, e que não são recuperáveis linearmente das representações de tokens isolados.

\paragraph{acl e advcl.} Com acurácia em torno de 89--93\%, as cláusulas adverbiais e
adnominais apresentam dificuldade moderada, refletindo a ambiguidade de adjunção em português:
uma mesma construção pode ser analisada como modificador nominal (\texttt{acl}) ou adverbial
(\texttt{advcl}) dependendo do contexto discursivo. Seus contornos de decisão envolvem
informação semântica que transcende o nível local da sentença.

\paragraph{nmod e obl.} Os rótulos nominais modificadores apresentam acurácia entre
93--96\%, mas são os que geram mais eventos de \textit{forgetting} durante o treinamento (ver
Seção~\ref{sec:instabilidade}), o que sugere que o modelo aprende padrões para esses rótulos,
mas os perde em ajustes posteriores de gradiente. A confusão frequente entre \texttt{nmod} e
\texttt{obl} — ambos expressando nominais associados a predicados ou outros nominais via
preposição, com distinção dependente de papéis temáticos — é coerente com a dificuldade de
codificação linear desse contraste (ver definições no Apêndice~\ref{cap:classe_deprel}).

A análise por rótulo revela ainda que a diferença entre modelos é mais pronunciada nos rótulos
de dificuldade moderada (90--95\%) do que nos extremos (rótulos fáceis $>$99\% ou muito
difíceis $<$50\%), onde a tarefa é inerentemente trivial ou insolúvel sem mais dados de
treinamento.

% ─────────────────────────────────────────────────────────────────────────────
\section{Análise de Instabilidade do Treinamento}
\label{sec:instabilidade}
% ─────────────────────────────────────────────────────────────────────────────

A análise de instabilidade, realizada com base nas predições por época sobre o conjunto de teste,
quantifica dois fenômenos distintos: \textbf{forgetting} (transição de acerto para erro em
épocas consecutivas) e \textbf{recovery} (transição inversa). Esses fenômenos indicam que o
modelo não converge monotonicamente: certos tokens são classificados corretamente em épocas
intermediárias mas perdem essa classificação em ajustes posteriores de gradiente.

\subsection{Curvas de Aprendizado do BERTimbau Large}

A Figura~\ref{fig:curvas_btl} mostra a progressão do UAS por época para as quatro variantes do
BERTimbau Large. Observa-se que o \textbf{cabeçote linear} converge de forma mais suave e
estável, enquanto o cabeçote \textbf{biaffine} apresenta oscilações mais pronunciadas nas
primeiras épocas, compatível com a maior instabilidade quantificada na análise de
\textit{forgetting}. As tabelas de treinamento completas por época — com UAS, LAS, UPOS,
\textit{training loss} e \textit{validation loss} para cada modelo — estão disponíveis no
Apêndice~\ref{cap:modelos_avaliados}, permitindo verificar em detalhe os pontos de inflexão
e as épocas de maior volatilidade de cada variante.

\begin{figure}[ht]
  \centering
  \includegraphics[width=0.92\textwidth]{Dissertação/Imagens/curvas_treinamento_btl.pdf}
  \caption{Curvas de UAS por época para as quatro variantes do BERTimbau Large. A variante
           Linear MTL converge de forma mais estável que a Biaffine MTL, que registra mais
           oscilações nas primeiras 20 épocas.}
  \label{fig:curvas_btl}
\end{figure}

Nota-se ainda que a variante Linear sem UPOS (sU) apresenta UAS final ligeiramente inferior ao
Linear MTL (+0,26 pp), confirmando o benefício marginal da tarefa auxiliar para o cabeçote
linear. Para o Biaffine, o comportamento é inverso: a variante sem UPOS alcança UAS maior que
a MTL, sugerindo que a tarefa de UPOS introduz conflito de gradientes nessa configuração.

\subsection{Instabilidade por Rótulo DEPREL}

A Figura~\ref{fig:instabilidade} apresenta o número de eventos de \textit{forgetting} para os
15 rótulos DEPREL mais instáveis do BERTimbau Large. Os resultados complementam a análise de
acurácia: enquanto rótulos como \texttt{parataxis} e \texttt{discourse} têm baixa acurácia por
dificuldade intrínseca, rótulos como \texttt{nmod} e \texttt{obl} têm acurácia razoável (93--96\%)
mas alta instabilidade, o que indica que o modelo os aprende e desaprende repetidamente durante
o treinamento. Isso sugere que as fronteiras de decisão para esses rótulos são próximas de
outros rótulos de alta frequência no espaço de representações, o que os torna vulneráveis a
deslocamentos de gradiente em épocas subsequentes.

\begin{figure}[ht]
  \centering
  \includegraphics[width=0.92\textwidth]{Dissertação/Imagens/instabilidade_deprel.pdf}
  \caption{Número de eventos de \textit{forgetting} (acerto\,$\to$\,erro) por rótulo DEPREL
           para as quatro variantes do BERTimbau Large. Rótulos de alta frequência com
           fronteiras de decisão difusas (\texttt{nmod}, \texttt{obl}) são os mais instáveis.}
  \label{fig:instabilidade}
\end{figure}

O rótulo \texttt{nmod} (modificador nominal) registra o maior número de \textit{forgettings}
em praticamente todas as variantes, com até 2.057 eventos na variante BtL-Biaff-MTL. Isso
se deve à alta ambiguidade de adjunção em português: sintagmas preposicionados podem ser
analisados como \texttt{nmod}, \texttt{obl} ou \texttt{advcl}, e a decisão correta depende
frequentemente de informação semântica não captada pelo cabeçote sintático. O rótulo \texttt{obl}
(oblíquo nominal) é o segundo mais instável, com padrão semelhante (para as definições formais
desses rótulos, ver Apêndice~\ref{cap:classe_deprel}).

A comparação entre variantes revela que a \textbf{variante Biaffine MTL é consistentemente a
mais instável}, com totais de \textit{forgetting} superiores às variantes Linear em todos os
rótulos analisados, exceto \texttt{conj} e \texttt{root}. Isso é coerente com a hipótese de
que o cabeçote biaffine, com maior capacidade e dependência de geometria das representações,
oscila mais ao longo do treinamento em corpus de tamanho limitado — e também com a hipótese da
codificação implícita: o cabeçote linear, ao operar como uma sonda (\textit{probe}) direta
sobre representações já organizadas sintaticamente pelo encoder, aprende uma fronteira de decisão
mais estável e menos suscetível a oscilações de gradiente.

A análise de instabilidade das Tabelas~\ref{tab:instab-mbert}--\ref{tab:instab-jab} no
Apêndice~\ref{ap:tabelas} mostra ainda que a diferença de instabilidade entre variantes com
e sem UPOS é pequena para o cabeçote linear (o MTL não aumenta nem reduz significativamente o
\textit{forgetting}), mas é mais expressiva para o biaffine (o MTL aumenta a instabilidade
em alguns rótulos).

% ─────────────────────────────────────────────────────────────────────────────
\section{Discussão Geral}
\label{sec:discussao_geral}
% ─────────────────────────────────────────────────────────────────────────────

Os resultados apresentados permitem extrair as seguintes conclusões:

\paragraph{1.\ A ausência do BiLSTM é o principal fator diferencial em relação ao PortParser.}
A lacuna de 1,09 pp em UAS em relação ao PortParser não é explicada pela diferença de
arquitetura do cabeçote, mas pela ausência da camada BiLSTM bidirecional que o PortParser
utiliza entre o encoder e o cabeçote biaffine. As curvas de treinamento no
Apêndice~\ref{cap:modelos_avaliados} mostram que os modelos convergem sem sinais de
sobreajuste grave, indicando que o teto de desempenho observado é imposto pela ausência da
contextualização adicional fornecida pelo BiLSTM, não pela capacidade de generalização do
cabeçote. Trabalhos futuros com a adição de um BiLSTM ao ParseH2IA devem reduzir
substancialmente essa lacuna.

\paragraph{2.\ O cabeçote linear é mais eficaz em corpus pequeno com ajuste fino completo,
possivelmente por explorar a codificação sintática implícita do BERT.}
A superioridade consistente do cabeçote linear sobre o biaffine em três dos quatro encoders,
mesmo em regime de ajuste fino completo, constitui o principal achado empírico deste trabalho.
A interpretação mais plausível, à luz dos resultados de instabilidade e da literatura de
\textit{probing}, combina dois mecanismos. Primeiro, estudos demonstram que modelos de
linguagem pré-treinados codificam implicitamente estrutura sintática em suas
representações~\citep{tenney2019bert,hewitt2019structural}: relações de dependência, distâncias
na árvore sintática e papéis gramaticais estão linearmente acessíveis no espaço vetorial
produzido pelo encoder. Um cabeçote linear funciona essencialmente como uma sonda direta
(\textit{linear probe}) sobre essa informação codificada, aprendendo de forma eficiente e
estável. O cabeçote biaffine, com suas quatro MLPs e operações bilineares, adiciona graus de
liberdade que não produzem ganho quando a informação relevante já está linearmente acessível.
Segundo, com apenas 5.893 sentenças de treino no Porttinari, a capacidade paramétrica
adicional do biaffine se manifesta como maior risco de sobreajuste em vez de melhor
generalização. A conclusão prática é que, ao trabalhar em corpus pequeno e sem BiLSTM,
arquiteturas de cabeçote mais simples não apenas bastam — elas são preferíveis. Essa conclusão
é reforçada pela consistência dos resultados da busca de hiperparâmetros
(Apêndice~\ref{cap:resultados_hiperparametros}): mesmo com configurações otimizadas
independentemente, o biaffine não alcança o linear em três dos quatro encoders.

\paragraph{3.\ O aprendizado multi-tarefa tem impacto marginal no \textit{parsing}.}
A tarefa auxiliar de UPOS, já próxima do teto de desempenho, não fornece gradiente suficiente
para melhorar substancialmente o desempenho em HEAD e DEPREL no Porttinari. O ganho marginalmente
positivo para os modelos lineares (+0,33 pp em UAS em média) e o impacto ligeiramente negativo
para os biaffines ($-$1,10 pp em média) sugerem que o sinal UPOS interage de forma distinta com
as duas arquiteturas de cabeçote: para o linear, atua como leve regularizador; para o biaffine,
pode introduzir conflito de gradiente entre as MLPs da tarefa principal e o classificador de UPOS.

\paragraph{4.\ JabuticaBERT apresenta desempenho abaixo do esperado.} A suposição de que um
modelo pré-treinado exclusivamente em português deveria superar o BERTimbau em tarefas para
o português não se confirma. Diferenças de tokenização e arquitetura do vocabulário podem
estar influenciando negativamente a extração de representações por subtoken nesta configuração.
Adicionalmente, a hipótese da codificação implícita sugere que a qualidade da codificação
sintática — e não apenas a língua do pré-treinamento — é determinante; nesse sentido, o
BERTimbau Large, por sua maior capacidade e treinamento, codificaria sintaxe de forma mais
rica que o JabuticaBERT Base. Nota-se também que os resultados da busca de hiperparâmetros
(Apêndice~\ref{cap:resultados_hiperparametros}) mostram maior variância (\textit{std LAS}) para
o JabuticaBERT em comparação aos demais encoders, o que indica maior sensibilidade aos
hiperparâmetros — possivelmente reflexo da mesma instabilidade de representações observada.

\paragraph{5.\ Rótulos de alta ambiguidade semântica permanecem desafiadores.} \texttt{parataxis},
\texttt{discourse} e \texttt{orphan} apresentam acurácia baixa e alta variabilidade entre
modelos, sugerindo que esses rótulos requerem informação discursiva e contextual que vai além
do que um cabeçote sintático treinado sobre o nível de sentença pode capturar — e que,
diferentemente dos rótulos gramaticalmente determinísticos, não está linearmente codificada
nas representações do encoder (ver definições formais no Apêndice~\ref{cap:classe_deprel}).

% ─────────────────────────────────────────────────────────────────────────────
\section{Limitações}
\label{sec:limitacoes}
% ─────────────────────────────────────────────────────────────────────────────

Este trabalho apresenta as seguintes limitações que devem ser consideradas na interpretação dos
resultados:

\begin{itemize}
  \item \textbf{Corpus único}: Os modelos foram treinados e avaliados exclusivamente no corpus
        Porttinari, o que restringe a generalização dos resultados a outros domínios e gêneros
        textuais. Avaliação em corpora de notícias, redes sociais ou textos jurídicos pode
        produzir resultados distintos.

  \item \textbf{Ausência de BiLSTM}: O ParseH2IA suprime a camada BiLSTM bidirecional presente
        no PortParser, o que impede comparação direta em condições arquiteturais idênticas. Os
        resultados refletem a escolha deliberada de avaliar o cabeçote de predição sem a influência
        da contextualização adicional fornecida pelo BiLSTM. Trabalhos futuros que adicionem o
        BiLSTM ao ParseH2IA permitirão isolar sua contribuição.

  \item \textbf{Ausência de decodificação estrutural}: Os modelos implementados utilizam
        decodificação gulosa (\textit{greedy}), sem garantir que a estrutura de árvore seja
        válida. O PortParser e outros modelos de estado da arte utilizam algoritmos de
        decodificação exata (e.g., algoritmo de Eisner) para garantir arborescências projetivas
        válidas, o que pode explicar parte da diferença de desempenho.

  \item \textbf{Avaliação de uma única execução}: Os resultados reportados correspondem a uma
        única execução de treinamento por configuração. Experimentos com múltiplas sementes
        aleatórias permitiriam estimar a variância dos resultados e tornar as comparações
        estatisticamente mais robustas.

  \item \textbf{Hiperparâmetros fixos por arquitetura}: A busca de hiperparâmetros foi conduzida
        individualmente para cada combinação de encoder e arquitetura
        (Apêndice~\ref{cap:resultados_hiperparametros}), replicando os intervalos de busca do
        PortParser original. Uma busca com espaço ampliado poderia beneficiar especialmente o
        cabeçote biaffine, que demonstrou maior sensibilidade aos hiperparâmetros.
\end{itemize}

Apesar dessas limitações, os resultados demonstram que é possível construir modelos de
\textit{parsing} competitivos para o português brasileiro com arquiteturas relativamente simples,
desde que se utilize um encoder pré-treinado de alta qualidade como o BERTimbau Large. A hipótese
de codificação sintática implícita, em particular, oferece uma contribuição interpretativa
relevante ao sugerir que a superioridade do cabeçote linear não é um artefato experimental,
mas reflexo de uma propriedade fundamental dos encoders BERT: a organização sintática latente
de suas representações.
